\subsection{Correct outcomes can still trigger justified rule warnings}
In the original test cohort, all 39 conversations with successful final outcomes also broke the interaction rule. Gemini warned in all 39. An evaluator that treated every successful outcome as error-free would therefore count every one of these warnings as a false alarm. The exact rule checker gives the same 39/39 warning count. There are no successful conversations that follow the rule in this cohort, so the corresponding warning rate is \emph{undefined}, not zero. Table~\ref{tab:rates} shows both denominators.

All 48 planned conversations completed. Gemini returned valid scores for 511 of 514 decisions. The three missing scores occur after an already observed first warning, so they do not prevent us from determining whether or when each conversation first triggered a warning. All 48 conversations violate the interaction rule, including the nine without state success. Table~\ref{tab:cohort} and Figure~\ref{fig:matrix} show the individual task and condition results.

\paragraph{Observed example: four lookups, one warning, a correct final state.}
In retail task 70 (DeepSeek actor, renamed tools), message 6 proposes four order lookups in one message. The applicable rule permits at most one tool call at a time. Gemini assigns score 1.0 and identifies the four simultaneous calls; this crosses the fixed 0.5 warning threshold. The conversation later ends with a matching final database. Thus the warning correctly detects this rule even though the final-outcome check passes. At the same recorded proposal, DeepSeek's monitor also describes the violation but assigns 0.35, below the warning threshold. Recognizing a rule in the explanation and issuing a thresholded warning are different observations. The calls are reads; this example establishes neither harmful execution nor advance prediction. The supplement's \texttt{release/examples.json} identifies the exact episode, zero-based message index, input hash and score records for this and the following examples.

\subsection{Gemini detects the visible proposal rather than warning in advance}
Gemini's first warning occurs at the first violating proposal in all 48 conversations: none occurs earlier or later, and no violating conversation lacks a warning. The monitor can already see the proposed calls at that point. This is detection of visible behavior, not evidence that it predicted an unseen failure. Figure~\ref{fig:timing} uses message order rather than elapsed time. The simple rule checker detects the same property directly; a language model is not necessary to recognize this particular violation.

The cohort contains 12 task IDs and 11 distinct reference-operation sequences. Both the task and operation-sequence bootstrap intervals for the fraction of state successes that violate the rule are [100.0\%, 100.0\%], with no undefined denominators in 5,000 resamples. Every observed state success violates the rule, so resampling these observations cannot produce a different fraction. These degenerate intervals provide no evidence that the population fraction is exactly one.
